\section{A quantitative full-entry correction margin}
\label{sec:terminal-modal-entry-correction-margin}

The directed decomposition in
\cref{lem:terminal-modal-directed-packet} uses only the sign $c\le0$.
At the last proper prefix the same correction ledger gives a little more.
The strengthening below is useful in the early half-retention interface, but
the resulting short-window convolution remains open.

We first sharpen the total-variation estimate only at the final prefix.  The
gain is small but uniform and leaves a useful rational margin.

\begin{lemma}[Sharpened final source variation; proved here]
\label{lem:terminal-modal-final-source-variation}
For the finite-$q$ derivative coefficients in
\eqref{eq:terminal-modal-finite-source-split},
\begin{equation}
 \operatorname{TV}(\varepsilon_n)<\frac1{300},
 \qquad
 |\varepsilon_{m-1}|+\operatorname{TV}(\varepsilon_n)
 <\frac{11}{1800}.
 \label{eq:terminal-modal-final-source-variation}
\end{equation}
Consequently the final derivative-packet loss in
\eqref{eq:terminal-modal-finite-derivative-loss} is below $11/1800$.
\end{lemma}

\begin{proof}
Retain $e_0(t)$, $A=\delta^{-r}$, $\lambda=-\log\delta/q$, and
$\kappa=-\log\chi/q$ from the proof of
\cref{lem:terminal-modal-finite-correction-sign}.  Put
\begin{equation}
 f(t)=e_0(t)-2te_0'(t)
 =\frac{t^4-30t^3+12t^2+10t+3}{10(1+t^2)^2}.
 \label{eq:terminal-modal-final-variation-f}
\end{equation}
The numerator of $f'$ after clearing its positive denominator is
\[
 p(t)=15t^4-10t^3-60t^2+6t+5.
\]
Now $p(7/8)=-137601/4096<0$, while
$p'(7/8)=-10467/128<0$ and
$p''(t)=60(3t+2)(t-1)\le0$.  Thus $f$ decreases on $[7/8,1]$.

Partition this interval into the sixteen cells
\[
 I_i=[7/8+i/128,\,7/8+(i+1)/128],\qquad0\le i<16.
\]
Monotonicity and exact rational summation give
\begin{equation}
 \sum_{i=0}^{15}\frac1{128\inf I_i}
 \max\{|f(\inf I_i)|,|f(\sup I_i)|\}<\frac{77}{12500}.
 \label{eq:terminal-modal-final-variation-cells}
\end{equation}
This is a finite rational certificate; the focused verifier performs the
sixteen evaluations and the exact comparison.

The parameter error relative to $f$ is at most
\begin{equation}
 (\lambda-1)|e_0|+(\kappa-2)|te_0'|+\frac{\lambda q}{10}
 \le\frac{2q}{5\delta}+\frac{2q^2}{5(1-q^2)}<\frac1{2500}.
 \label{eq:terminal-modal-final-variation-parameter}
\end{equation}
Under the change of variables $t=\chi^{r+1}$, one has
$|dr|=|dt|/(\kappa qt)$.  Since $A\le16/15$, $\kappa\ge2$, and
$\log(8/7)<1/7$, the parameter-error integral inside the common
$8/15$ factor is strictly below
$(1/2500)\log(8/7)<1/17500$.  Equations
\eqref{eq:terminal-modal-final-variation-cells}--
\eqref{eq:terminal-modal-final-variation-parameter} yield
\begin{equation}
 \operatorname{TV}(E_n)
 <\frac8{15}\left(\frac{77}{12500}+\frac1{17500}\right)
 =\frac{1088}{328125}<\frac1{300}.
\end{equation}
Finally $\operatorname{TV}(\varepsilon_n)=\operatorname{TV}(E_n)$ and
$|\varepsilon_{m-1}|\le1/360$ by
\eqref{eq:terminal-modal-finite-epsilon-pointwise}.  Adding the two bounds
gives \eqref{eq:terminal-modal-final-source-variation}.  The Abel multiplier
four cancels the definition $h_n=\varepsilon_n/4$, exactly as in
\eqref{eq:terminal-modal-finite-derivative-loss}.
\end{proof}

\begin{lemma}[Sharpened final-prefix correction; proved here]
\label{lem:terminal-modal-sharpened-final-prefix-correction}
Let $m\ge64$.  On the coordinates shared by the last two proper prefixes,
\begin{equation}
 \frac{K_m(j)}{q^3(1-q)^{m-2}}
 \ge \frac{16649}{230400}>\frac7{100},
 \qquad 0\le j\le m-2.
 \label{eq:terminal-modal-sharpened-final-prefix-shared}
\end{equation}
At the moving frontier,
\begin{equation}
 \frac{K_m(m-1)}{q^3(1-q)^{m-2}}>\frac25.
 \label{eq:terminal-modal-sharpened-final-prefix-frontier}
\end{equation}
\end{lemma}

\begin{proof}
Put $\delta=1-q$.  We first sharpen only the final-time leading margin in
the proof of \cref{lem:terminal-modal-finite-correction-sign}.  Retain the
notation $d_n^\circ(j)$, $p_k$, and $h_{j,k}$ from
\eqref{eq:terminal-modal-leading-correction-fixed-j}, so that, for $j\ge2$
and $k=n-j-1$,
\begin{equation}
 d_n^\circ(j)=\frac{p_k-2h_{j,k}}{40},
 \qquad p_k=8+4(-1/2)^k\quad(k\ge2).
 \label{eq:terminal-modal-final-leading-coefficient}
\end{equation}
The two exceptional seed rows are already explicit in the proof of
\cref{lem:terminal-modal-leading-correction-sign}.  For $n=m\ge64$ they give
\begin{equation}
 d_m^\circ(0)\ge\frac1{10}-\frac1{10\cdot32}
 =\frac{31}{320},\qquad d_m^\circ(1)\ge\frac1{10}.
 \label{eq:terminal-modal-final-leading-seed-neighbor}
\end{equation}

For completeness, here are all edge cases needed to sharpen the remaining
rows.  The recurrence
\begin{equation}
 h_{j+1,k}=\frac{h_{j,k}+h_{j+1,k-1}}2,
 \qquad h_{j,-1}=0,
 \label{eq:terminal-modal-final-h-recurrence}
\end{equation}
and the displayed formula for $h_{2,k}$ in the earlier proof give
\begin{align}
 h_{3,0}&=\frac94,&h_{3,1}&=\frac52,&
 h_{3,k}&\le2+2^{-k}\quad(k\ge2),
 \label{eq:terminal-modal-final-h3}\\
 h_{4,0}&=\frac98,&h_{4,1}&=\frac{29}{16},&
 h_{4,2}&=\frac{65}{32},&h_{4,3}&=\frac{131}{64}.
 \label{eq:terminal-modal-final-h4-edge}
\end{align}
A partial-fraction extraction gives, for every $k\ge0$,
\begin{equation}
 h_{4,k}-2=
 \frac{-(-1)^k-129-66(k+1)
 +124\binom{k+2}{2}-40\binom{k+3}{3}}{2^{k+7}}.
 \label{eq:terminal-modal-final-h4-formula}
\end{equation}
The numerator is zero at $k=4$.  Its polynomial part has forward difference
$62+24k-20k^2$, while the alternating term changes that difference by at
most two.  Hence it strictly decreases after $k=4$, and
$h_{4,k}\le2$ for $k\ge4$.

Applying \eqref{eq:terminal-modal-final-h-recurrence} once more gives
\begin{equation}
 h_{5,0}=\frac9{16},\quad h_{5,1}=\frac{19}{16},\quad
 h_{5,2}=\frac{103}{64},\quad h_{5,3}=\frac{117}{64}.
\end{equation}
Thus $h_{5,k}\le2$ follows by induction on $k$, using the $j=4$ bound from
$k=4$ onward.  A double induction in $j\ge5$ and $k\ge0$ in
\eqref{eq:terminal-modal-final-h-recurrence} now proves
\begin{equation}
 h_{j,k}\le2\qquad(j\ge5, k\ge0).
 \label{eq:terminal-modal-final-h-double-induction}
\end{equation}
The bases are $h_{j,0}=9/2^{j-1}\le9/16$, and the induction step averages
two quantities already at most two.

We apply these bounds only on the final diagonal $j+k=m-1$.  If $k\ge4$,
then $p_k\ge63/8$ except for even $k$, where $p_k>8$; hence
\eqref{eq:terminal-modal-final-leading-coefficient} and
\eqref{eq:terminal-modal-final-h-double-induction} give $31/320$.  The row
$j=4$ uses \eqref{eq:terminal-modal-final-h4-formula}; the row $j=3$ has
$k\ge60$, and \eqref{eq:terminal-modal-final-h3} gives the same strict
bound because $6/2^{60}<1/8$; the row $j=2$ uses $h_{2,k}\le2$ for
$k\ge2$.  For the three frontier diagonals, $k=1,2,3$, one has
$j\ge60$.  At $k=1$,
\begin{equation}
 h_{j,1}=\frac{9j-7}{2^j}\le\frac{47}{64}<\frac{17}{16}
 \qquad(j\ge6),
\end{equation}
while $k=2$ follows from $h_{j,2}\le2$.  Finally, from the defining product
for $h_{j,k}$ and positivity of the coefficients of
$((1-t)(1+t/2))^{-1}$,
\begin{equation}
 h_{j,3}
 \le2^{-j}\left\{18\sum_{r=0}^3
 \frac{\binom{j+r-1}{r}}{2^r}+5\right\}
 \le2^{-j}\{72(j+2)^3+5\}<\frac1{16}
 \quad(j\ge60).
\end{equation}
Here the last majorant decreases with $j$, and its value at $60$ is less
than $1/16$.  Substitution for $k=1,2,3$ proves in every case
\begin{equation}
 d_m^\circ(j)\ge\frac{31}{320},\qquad0\le j\le m-2.
 \label{eq:terminal-modal-final-leading-margin}
\end{equation}

We next repeat the finite-$q$ perturbation ledger only at final time $N=m$.
The sharpened stopped derivative loss is $11/1800$ by
\cref{lem:terminal-modal-final-source-variation}, and the initial-pair loss
is at most
\begin{equation}
 \frac{3q}{5}\le\frac3{5120}.
\end{equation}
The point-mass loss improves by a factor two at this final time.  Indeed, a
source born at $n$ is propagated for $k=m-n$ steps.  For $k\ge2$ the line
coefficient $\ell_k(r)$ is at most $1/2$ and vanishes unless
$|r|\le k-1$.  At a final shared coordinate $0\le j\le m-2$, the positive
and reflected source branches can be nonzero only if, respectively,
\begin{equation}
 2n\le m+j-1,qquad 2n\le m-j-1.
 \label{eq:terminal-modal-final-mass-support}
\end{equation}
There is no second cycle alias because $k<m$, and the $k=1$ source at
$n=m-1$ misses every shared coordinate.  The sum of the two support counts
is at most $m-1$.  Since each nonzero coefficient is at most $1/2$ and
$|\mu_n|<43q/75$, the total mass loss is strictly less than
\begin{equation}
 \frac{43q(m-1)}{150}<\frac{43}{2400}.
 \label{eq:terminal-modal-final-mass-loss}
\end{equation}
This count includes the reflected branch and the $k=1$ edge case; no
infinite-line mass estimate is substituted for the folded response.

Combining \eqref{eq:terminal-modal-final-leading-margin}, the derivative,
base, and mass bounds gives
\begin{equation}
 \frac{K_m(j)}{q^3\delta^{m-2}}
 \ge\frac{31}{320}-\frac{11}{1800}-\frac3{5120}-\frac{43}{2400}
 =\frac{16649}{230400},
\end{equation}
which proves \eqref{eq:terminal-modal-sharpened-final-prefix-shared}.

It remains to treat the moving frontier.  Put
$x_n=\bar c_n(n-1)$ in the rescaling
\eqref{eq:terminal-modal-finite-rescaling}.  The exact new-row recurrence is
\begin{equation}
 x_{n+1}=\frac12x_n+F_n,qquad
 x_2=-\frac q4+\frac{q^2}{20}>-\frac q4.
\end{equation}
The source split gives $F_n=9/40-3\varepsilon_n/4+\mu_n$.  Therefore
\begin{equation}
 F_n>\frac9{40}-\left(\frac9{40}+\frac{43}{75}\right)q
 =\frac9{40}-\frac{479}{600}q>\frac7{32},
\end{equation}
where the final slack is $3361/614400$ at $q=1/1024$.
Iteration and $2^{-(m-2)}\le1/16$ give
\begin{equation}
 x_m>\frac7{16}(1-2^{-(m-2)})
       -\frac{q}{4\,2^{m-2}}
 \ge\frac{105}{256}-\frac1{65536}>\frac25.
\end{equation}
Since the preceding correction is zero-padded at the new coordinate,
$x_m=K_m(m-1)/(q^3\delta^{m-2})$.  This proves
\eqref{eq:terminal-modal-sharpened-final-prefix-frontier}.
\end{proof}

\begin{lemma}[Uniform negative entry-correction margin; proved here]
\label{lem:terminal-modal-entry-correction-margin}
For every $m\ge64$, the literal full-face entry correction in
\eqref{eq:terminal-modal-propagator-data} satisfies
\begin{equation}
 c(j)<-\frac{13}{100}q^3,\qquad0\le j\le m.
 \label{eq:terminal-modal-entry-correction-margin}
\end{equation}
\end{lemma}

\begin{proof}
Let $L_m=(I+P_m^{\rm row})/2$ be the lazy substochastic path stencil on the
last proper prefix, and let $e_{m-1}$ denote its last coordinate vector.
Define the positive frontier scalar
\begin{equation}
 a_m=\frac{q\eta P_{m-1}}2-\frac{q^3}{5}>0,qquad
 \frac{a_m}{q^3}\le\frac{33}{64}-\frac15=\frac{101}{320}.
 \label{eq:terminal-modal-entry-frontier-defect-scalar}
\end{equation}
The upper bound uses $P_{m-1}/q^2\le33/16$ and $\eta\le1/2$.

The exact old-row correction identity, including its frontier defect, is
\begin{equation}
 c|_{[0,m-1]}
 =-2\delta L_mK_m+\delta a_mL_me_{m-1}+g_m e_0.
 \label{eq:terminal-modal-entry-two-row-defect}
\end{equation}
To verify it, write the correction part of the final average negative
residual as
\[
 Z_m^C=\frac{2K_m-a_me_{m-1}}{1+q}.
\]
The post-step residual is $B_mR_m$, and
$B_m/(1+q)=\delta L_m$.  The ideal post-step residual equals $g_j$ for
$1\le j<m$ and equals $g_0+g_m$ at the reflected seed.  Subtraction gives
\eqref{eq:terminal-modal-entry-two-row-defect}.  In particular the positive
defect is supported on exactly two rows:
\begin{equation}
 (L_me_{m-1})(m-2)=\frac14,qquad
 (L_me_{m-1})(m-1)=\frac12,qquad
 (L_me_{m-1})(j)=0\ (j\le m-3).
 \label{eq:terminal-modal-entry-two-row-support}
\end{equation}
This is the corrected identity; omitting either row is false.

Put $A=16649/230400$.  For $j\le m-3$, positivity of the lazy stencil,
\cref{lem:terminal-modal-sharpened-final-prefix-correction}, and
$\delta^{m-1}\ge1-(m-1)q>15/16$ give
\begin{equation}
 -\frac{c(j)}{q^3}
 \ge2A\delta^{m-1}
 >\frac{16649}{122880}>\frac{13}{100};
 \label{eq:terminal-modal-entry-interior-margin}
\end{equation}
the additional seed term $g_m e_0\le0$ only helps.  At the two frontier
rows, \eqref{eq:terminal-modal-entry-two-row-support} and the shared/frontier
bounds give respectively
\begin{align}
 -\frac{c(m-2)}{q^3}
 &>2\frac{15}{16}\left(\frac34A+\frac1{10}\right)
     -\frac{101}{1280}
 =\frac{34441}{163840}>\frac{13}{100},
 \label{eq:terminal-modal-entry-penultimate-margin}\\
 -\frac{c(m-1)}{q^3}
 &>2\frac{15}{16}\left(\frac15+\frac14A\right)
     -\frac{101}{640}
 =\frac{123401}{491520}>\frac{13}{100}.
 \label{eq:terminal-modal-entry-frontier-margin}
\end{align}

Finally the actual full-face endpoint is not covered by the degree-two
prefix stencil.  Its exact value is
\begin{equation}
 c(m)=\frac{q^3}{5}-\eta\sigma_m-g_m.
\end{equation}
The raw-frontier estimate already used in the proof of
\cref{lem:terminal-modal-directed-packet} gives
\[
 \frac{\eta\sigma_m-q^3/5}{q^3}
 \ge\frac{6985811}{9835520},
 \qquad
 \frac{|g_m|}{q^3}=\frac{8m\delta^m}{2^m}<\frac18.
\]
Therefore
\begin{equation}
 -\frac{c(m)}{q^3}>
 \frac{6985811}{9835520}-\frac18
 =\frac{5756371}{9835520}>\frac{13}{100}.
\end{equation}
This completes \eqref{eq:terminal-modal-entry-correction-margin}.
\end{proof}

The sole early nonlinear target remains the correlated inequality
\eqref{eq:terminal-modal-early-half-kernel}.  The new theorem improves its
$-\mathcal H_kc$ term, but neither discards nor separately upper-bounds
$\mathcal L_ku$.

Together with the sharpened position stitch, the exact remaining time window
can nevertheless be shortened.  Since $\mathcal H_k$ is positive with row
mass $1/2$, \cref{lem:terminal-modal-entry-correction-margin} gives
\begin{equation}
 -\mathcal H_kc>\frac{13q^3}{200}.
\end{equation}
Consequently the finite inequality
\begin{equation}
 \mathcal L_ku\le\frac{13q^3}{200}
 \qquad(1\le k,\ qk<3/50)
 \label{eq:terminal-modal-short-early-convolution-open}
\end{equation}
would imply strict half-retention throughout that shorter interval; the
directed ideal contribution in
\eqref{eq:terminal-modal-early-half-kernel} is already nonnegative.  The late
position comparison begins at $qk=3/50$.  Inequality
\eqref{eq:terminal-modal-short-early-convolution-open} remains \textbf{OPEN};
the analogous claim through $17/200$ is false and is not used.
